% Transcription — fonds Alexandre Grothendieck, shelfmark 75, batch 1 (pages 1-20).
% Established from the facsimile archives/batches/75/batch-01.pdf (prepared
% from a copy of 75.pdf uploaded by the user; PDF page k = archivists' page k,
% no cover sheet), per the skill transcribe-grothendieck. The folder is
% thirty-two pages; this is the first of two batches. Pages 21-32 are batch 2,
% transcribed in parallel by another pass; its notation was read and followed
% where the two meet (omega for the two-element set of orientations,
% G+ = Aut+, the contracted product written with a wedge and the group as
% subscript, \mathfrak{S}_n / \mathfrak{A}_n, \mathbf{Z}, \mathbf{R}).
%
% The dating is the folder's own, copied verbatim from the catalogue; the
% group « Géométrie et topologie combinatoire » (69-89, 1976-[vers 1986])
% holds it, and since the folder has a date of its own the group's range is
% not added. Folder 75 is « Cours C4, DEA : notes manuscrites (s.d.), lettre
% (1978), tapuscrit (1977-1978) ». The letter and the typescript are both in
% batch 2 (pp. 28-32); nothing in this batch is typed, and nothing is by
% another hand.
%
% Pages skipped: none. Page 1 is the folder's cover, inscribed in pencil along
% the right edge « C4 DEA — Géométrie combinatoire », with « (24 p) » above;
% by the rule for covers it takes no \page marker, and its words become the
% first section heading, with a \note.
%
% Pages summarised: none. Pages transcribed from his hand: 2-20, black ink
% throughout, with a few pencil additions on p. 4 and p. 5. The hand is fast
% course-preparation drafting: the formulas carry, the connecting prose
% mostly carries, and the densest lines (the ends of pp. 13, 15, 18, 20) do
% not, and are left \ill{}.
%
% His pagination: none on pp. 2-15. From p. 16 he numbers the items of a run
% with circled numerals: ① (« Yoga », p. 16), ②, ③, ④ (p. 17), ⑤ (p. 19),
% ⑥, ⑦ (p. 20). Batch 2 opens on p. 21 with a circled « 8 » (« Digression
% sur structures cubiques »), so the run continues across the batch boundary
% and pp. 16-20 hold its items 1-7. Item ⑦ on p. 20 was begun (« Orientations
% ... X ∧ Y ») and cancelled; « Orientations » is taken up again on p. 24.
% Earlier, pp. 8-12 carry their own numbered sequence 1)-4) (1) sets with a
% permutation, 2) « Rectification », 3) semi-direct products, 4) sets of small
% cardinal, a)-d)).
%
% What the batch is about. Preparation for the 1977-78 C4 and DEA courses
% on the combinatorial geometry of polygons, built on torsors. P. 2: groups
% acting on sets, homogeneous spaces as subgroups up to conjugacy, torsors,
% Aut_G P ≃ G non-canonically with u_{sx} = u_x ∘ int(s^{-1}). Pp. 3-5
% (« Graphes »): a graph as (S, A, R), strict graphs, the geometric
% realisation by segments, a Proposition that a finite graph has a rectilinear
% realisation iff it is strict, and then one in R^3; contours (every vertex on
% two edges), classification of connected contours by card S (subgroups of Z),
% orientation automorphisms, ω(X) of cardinal 2, Aut+(X) cyclic of index 2,
% Aut(X) ≃ Z/2Z.Z/H. Pp. 6-7: an oriented contour is a torsor under
% Z/nZ; the data (cyclic group G^0, pair {g, g^{-1}} of generators, G^0-torsor
% S); orientation as a choice of origin σ(a) of each edge, u_ω, u_{-ω} =
% u_ω^{-1}, S and A as torsors under G^0. Pp. 8-9: sets with permutation =
% Z-sets, conjugacy classes in S_n and partitions (a small table for S_2-S_4,
% A_2-A_4), the « Rectification » that oriented contours are permutations
% without fixed point, semi-direct products and the affine group
% Aff(Z) ≃ Z.Aut(Z), Aff_R(V), displacements V.Orth(V, Q). Pp. 10-11: the
% sections q_t of 1 -> Z -> Aut(Z, T) -> Aut(Z) -> 1 and q_{tg} = int(i(g)) q_t;
% the dihedral group D_n = Z_n.{±1}; σ_{tg} = g² σ_t, so for n odd the vertex
% t -> σ_t is a bijection onto the non-oriented automorphisms, fibres of
% cardinal 2 for n even (opposite vertex); fibre products. Pp. 12-13:
% pull-back of extensions; small symmetric groups as affine/projective groups
% (S_2 ≃ Aff(1,F_2), S_3 ≃ Aff(1,F_3) ≃ Gl(2,F_2) ≃ GP(1,F_2), S_4 ≃
% Aff(2,F_2) ≃ GP(1,F_3), A_4 ≃ Aff(1,F_4) ?, S_5 ≃ GP(1,F_5), A_5 ≃
% GP(1,F_4)); vector, affine and projective structures as elements of
% Isom(k^n, V)/Gl(n,k), .../Aff(n,k), Isom(P^{n-1}(k), P)/GP(n-1,k). P. 14:
% exercise on regular realisations of polygons in the Euclidean plane
% (existence by z^n = 1, uniqueness up to similitude), affine hull, centre of
% gravity. Pp. 15-16: regular polygons again; functoriality of P ×_G E = ^P E,
% P -> Bij(E, F) injective when G acts faithfully, transport of structure,
% the « Yoga »: objects isomorphic to E <-> torsors under Aut(E),
% F -> Isom(E, F), P -> P ×_G E. Pp. 17-18: ② G-torsors as (G, E)-structures,
% ③ as structures of species C isomorphic to E, ④ the vector hull of a
% combinatorial polygon via the regular realisation ζ^i in C and D_n acting on
% C, Env(P) = ^T E with T = Isom(P_n, P); Hom(Tors(G), C) ≃ G-objects of C,
% C_E ≃ Tors(G), Hom(C_E, E) ≃ G-(E) (marked « DEA » in the margin).
% Pp. 19-20: ⑤ extension of the structure group, E ×_G G', twisting by
% ^{G'}T; ⑥ the product of torsors under a commutative group, X ∧ Y ∧ Z, the
% product of torsors is a torsor, inverses Y ≃ X ×_G (G --(g -> -g)--> G);
% ⑦ cancelled.
%
% Notation fixed once for the batch: his gothic S and A as \mathfrak{S}_n,
% \mathfrak{A}_n (as batch 2); double-struck letters as \mathbf{Z}, \mathbf{R},
% \mathbf{C}, \mathbf{F}_q, \mathbf{P}; his « Gl », « GP », « Aff », « Aut »
% kept; the contracted product written as he writes it on each line, × or ∧
% with the group as subscript (he uses both, sometimes on the same page, and
% the choice is not regularised); twisting as ^T E (left superscript); the
% underlined category names as \underline{\mathrm{Hom}}, \underline{\mathrm{Tors}};
% « ½ direct » kept for semi-direct. His circled item numbers as
% (n). His \boxed formulas kept boxed.
%
% Differences with batch 2 worth noting: none contradict it. Batch 2's
% header says the circled « 8 » of p. 21 presumably continues a run begun in
% batch 1; it does — items ① to ⑦ are on pp. 16-20 here.
%
% Readings to check first: the cover inscription (« Géométrie
% combinatoire »); « non » at the end of the connectedness line on p. 4 and
% « Dénombrement » there; the words of the item-6 fragment on p. 6 (« choix »,
% « non connexes »); the partition table on p. 8, especially the cancelled
% A_4 line, and the small oblique margin note at its foot (H², extensions);
% « non dégénérée » on p. 9; the diagonal margin note on p. 11 (opposite
% vertices); the bracketed paragraph on abstract affine spaces on p. 13; the
% phrase between « et tous les angles itou » and « unimodulaire » on p. 14;
% the bottom of p. 15 (the long bracket on (G, E)-structures); « liée » and
% the whole lower half of p. 18 (the Hom(C_E, E) square and the diagonal
% margin note); the opening of p. 19 and the twisting sentence that runs from
% p. 19 onto p. 20.
%
% Pass: Opus 5.5 (claude-opus-5-5), 2026-09-24 — first pass, unchecked against the pages by a human.

\documentclass[11pt,a4paper]{article}
% Le préambule commun à toutes les transcriptions du fonds Grothendieck.
%
% Il tient en une idée : **l'incertitude se compose, elle ne se résout pas.**
% Une transcription qui lisse un mot douteux a détruit la seule chose pour
% laquelle on la faisait — un lecteur doit pouvoir savoir, à chaque endroit,
% ce qui était sur le feuillet et ce qui a été ajouté. Les six macros qui
% suivent sont donc l'appareil critique, et non de la décoration.
%
% Les mêmes commandes sont reconnues par `scripts/render.mjs`, qui produit la
% vue de lecture du site. Une macro ajoutée ici doit l'être aussi là-bas,
% faute de quoi le rendu échouera bruyamment — ce qui est voulu : un rendu
% silencieusement faux est pire qu'un rendu absent.

\ProvidesPackage{grothendieck}[2026/08/08 Transcriptions du fonds Grothendieck]

\usepackage{amsmath,amssymb,amsthm}
\usepackage{mathrsfs}
\usepackage{stmaryrd}
% Les diagrammes commutatifs sont partout dans ce fonds : c'est la langue
% même de la géométrie algébrique de Grothendieck, et une transcription qui
% les rendrait en prose ne transcrirait rien.
\usepackage{tikz-cd}
% Français par défaut — c'est la langue des feuillets — mais l'anglais est
% chargé aussi : la traduction et le résumé sont des documents anglais, et la
% typographie française (espaces avant les deux-points) y serait une faute.
% Ces documents basculent par \selectlanguage{english} en tête de corps.
\usepackage[english,french]{babel}
\usepackage{fontspec}
\usepackage{geometry}
\geometry{a4paper,margin=25mm,marginparwidth=28mm}
\usepackage{marginnote}
\usepackage{xcolor}
% ulem, not soul. `soul`'s \st and \ul are built on a character-by-character
% scan that XeLaTeX's font machinery breaks: the document fails with an
% undefined control sequence rather than degrading. ulem does the same two jobs
% and is Unicode-engine safe. `normalem` stops it hijacking \emph, which the
% transcriptions use for Grothendieck's own emphasis.
\usepackage[normalem]{ulem}
\usepackage{hyperref}
\hypersetup{colorlinks=true,linkcolor=black,urlcolor=black,pdfborder={0 0 0}}

% --- Métadonnées du lot -----------------------------------------------------
% Reprises telles quelles par le script de rendu, qui les affiche en tête de
% la vue de lecture. Elles fixent ce que le fichier prétend couvrir : sans
% elles, une transcription flotte sans point d'attache dans un fonds de seize
% mille feuillets.

\newcommand{\folder}[1]{\gdef\grothFolder{#1}}
\newcommand{\batch}[1]{\gdef\grothBatch{#1}}
\newcommand{\leaves}[2]{\gdef\grothFirst{#1}\gdef\grothLast{#2}}
\newcommand{\foldertitle}[1]{\gdef\grothFoldertitle{#1}}
\gdef\grothFolder{?}\gdef\grothBatch{?}\gdef\grothFirst{?}\gdef\grothLast{?}\gdef\grothFoldertitle{}

% --- Le résumé --------------------------------------------------------------
%
% Il ouvre la lecture modernisée, sous le titre « Résumé », et il est composé
% en petit. Ce n'est pas un ornement : c'est le seul passage du document qui
% ne soit pas des mathématiques, et un lecteur doit pouvoir le reconnaître
% avant de l'avoir lu — puis le sauter s'il connaît déjà le sujet, ou s'y
% arrêter s'il ne le connaît pas. Le filet qui le suit marque la reprise du
% corps.

\newenvironment{resume}
  {\small\setlength{\parskip}{0.45em}}
  {\par\medskip\hrule\medskip}

% --- Mots-clés ---------------------------------------------------------------
%
% En anglais, en clôture du résumé de la lecture modernisée : le vocabulaire
% moderne sous lequel chercher ce que les feuillets construisent. Le manifeste
% du site (`npm run manifest`) les extrait pour étiqueter la cote entière —
% cette ligne est la source unique des tags, il n'y a pas de fichier à côté.
\newcommand{\keywords}[1]{\par\smallskip\noindent
  {\footnotesize\textsc{Keywords} --- \textit{#1}}\par}

% --- L'appareil critique ----------------------------------------------------

% Le numéro de feuillet, en marge. C'est l'ancre qui permet de régler un
% désaccord sur une formule en regardant *un* feuillet plutôt qu'une chemise
% de six cents. Le site s'en sert aussi pour tourner les pages du fac-similé
% quand on fait défiler la transcription.
\newcommand{\leaf}[1]{%
  \marginnote{\footnotesize\color{gray!70}\textsf{#1}}%
  \label{leaf:#1}\ignorespaces}

% Illisible : on ne devine pas. Un mot inventé qui se lit comme les autres est
% le pire résultat possible.
\newcommand{\ill}{\textcolor{red!60!black}{[\,\ldots\,]}}

% Lecture douteuse : le mot est donné, et signalé comme incertain.
\newcommand{\uncertain}[1]{\uline{#1}}

% Ajout de l'éditeur — une lettre manquante, un mot sous-entendu.
\newcommand{\add}[1]{\textcolor{blue!50!black}{[#1]}}

% Biffé par Grothendieck lui-même. Ce qu'il a rayé fait partie du document :
% une rature dit souvent où la pensée a changé de direction.
\newcommand{\struck}[1]{\sout{#1}}

% Note du transcripteur, sur l'état matériel du feuillet.
\newcommand{\note}[1]{\par\smallskip{\small\color{gray!60!black}\textit{[#1]}}\par\smallskip}

% Note marginale de Grothendieck, distincte du corps du texte.
\newcommand{\marginal}[1]{\par\smallskip\noindent\hspace{1em}%
  {\small\color{gray!60!black}#1}\par\smallskip}

% --- Titre ------------------------------------------------------------------

\AtBeginDocument{%
  \begin{center}
    {\large\bfseries Fonds Alexandre Grothendieck --- cote n\textsuperscript{o}~\grothFolder}\\[2pt]
    {\itshape\grothFoldertitle}\\[4pt]
    {\small feuillets \grothFirst--\grothLast\ (lot \grothBatch)}\\[2pt]
    {\footnotesize\color{gray!60!black}Transcription établie d'après le fac-similé numérisé
     par l'Université de Montpellier.\\
     Les lectures incertaines sont soulignées, les illisibles notées [\,\ldots\,],
     les ajouts entre crochets.}
  \end{center}
  \vspace{1em}}

\endinput


\folder{75}
\batch{1}
\pages{1}{20}
\dating{1977-1978}
\watermark{Édition de démonstration}
\foldertitle{Cours C4, DEA : notes manuscrites (s.d.), lettre (1978), tapuscrit (1977-1978)}

\begin{document}

\section*{C4 DEA — Géométrie combinatoire}

\note{inscrit au crayon, le long du bord droit, sur la chemise (page 1), qui
ne porte rien d'autre ; au-dessus, entre parenthèses, « (24 p) ». La chemise
ne reçoit pas de numéro de page}

\section*{Rappels de théorie des groupes et sur espaces homogènes}

\page{2}
Groupe (loi comp.\ assoc.\ unitaire avec inverses)

hom de groupes (on a néc.\ $f(e) = e$, $f(x^{-1}) = f(x)^{-1}$)

opération d'un groupe $G$, sur un ens (ou : groupe d'opérateurs)

opération transitive ($\forall x, y \in E$, $\exists g \in G$, $y = gx$) —
espace homogène (on exige de plus $E \neq \emptyset$)

(couples $(E, x)$, $E$ esp.\ hom., $x \in E$) $\Longleftrightarrow$ donnée d'un
ss-gpe $H$ de $G$

c'est en fait une « équivalence de catégories »

\begin{tikzcd}[row sep=small]
  (E, x) \arrow[d, "f"'] & & \\
  (E', x') & (f(x) = x') & \Longleftrightarrow \; H \subset H'
\end{tikzcd}

\emph{NB} $f$ (commutant à $G$) connu quand on connaît $f(x)$, si $G$
transitif sur $E$.

$E \sim E'$ ssi $H$, $H'$ conjugués ; $\exists\, E \to E'$ ssi $H$ conjugué à un
ss-groupe de $H'$

\marginal{classes d'iso.\ d'esp.\ hom.\ sous $G$ \uncertain{$\Leftrightarrow$}
classes de \uncertain{conjugaison} de s-groupes}\note{écrit en diagonale dans
la marge gauche, à hauteur du diagramme}

\emph{Torseur sous $G$} :
\[
  \left\{
  \begin{array}{l}
    \forall x, y \in E,\ \exists !\, g \in G,\ gx = y \\
    E \neq \emptyset
  \end{array}
  \right\}
  \Longleftrightarrow
  \left\{
  \begin{array}{l}
    E \text{ esp. hom., stabilisateurs} \\
    \text{de pts réd. à } \{e\}
  \end{array}
  \right\}
  \Longleftrightarrow
  \{ E \simeq G_s \}
\]

Les torseurs sous $G$ sont isomorphes (non canoniquement).

Les torseurs sous $G$ sont les « plus gros » espaces homogènes sous $G$.

[le groupe des autom.\ d'un torseur $P$ sous $G$ est $\simeq G$ (non
canoniquement) : le choix de $x \in P$ définit un iso.\
$G \xrightarrow[\sim]{u_x} \operatorname{Aut}_G P$, si on remplace $x$ par
$y = sx$, cet iso.\ est remplacé par
\[
  \boxed{u_y = u_{sx} = u_x \circ \operatorname{int}(s^{-1})}
\]
\note{dans le cadre, un premier essai après le deuxième signe $=$ est biffé}
\begin{gather*}
  u_x(g).x = gx \\
  u_y(g)\,y = gy = gsx \\
  u_y(g)\,sx = s\,u_y(g)\,x \quad \text{?} \\
  u_y(g).x = s^{-1}gsx = u_x(s^{-1}gs).x
\end{gather*}
\struck{$(u_x \circ \operatorname{int}(s)^{-1})(g)\,y = u_x(sgs^{-1})(y) =$}
\[
  u_y(g) = u_x(s^{-1}gs) \quad ]
\]

\section*{Graphes}

\page{3}
\note{au-dessus du titre, un mot biffé, illisible : \struck{\ill{}}}

\emph{Déf} \emph{Graphes}
\[
  \left.
  \begin{array}{l}
    \text{Ens } S \text{ de « sommets »} \\
    \text{ens } A \text{ « d'arêtes »} \\
    R \subset S \times A \quad \text{relation d'incidence}
  \end{array}
  \right\} \text{Données}
\]

\emph{Axiomes} 1) $\forall a \in A$, $\operatorname{card} R\{a\} = 2$
(chaque arête a exactement $2$ sommets)

\note{à gauche, un dessin de graphe : une douzaine de sommets marqués d'un
point, des arêtes droites et courbes, deux arêtes doubles joignant les mêmes
sommets, une arête isolée et un sommet isolé}

Un graphe est dit \struck{\ill{}} \emph{strict} si une arête est connue
quand on connaît ses sommets.

\struck{Une} La donnée d'un \emph{graphe strict} équivaut donc à la
donnée d'un ens.\ $S$, et d'une partie $A$ de $\mathfrak{P}_2(S)$ (ens.\ des
parties de card.\ $2$ de $S$).

Arc d'un graphe (ou arête orientée) — ens.\ $\vec{A}$

\struck{Remarque} \emph{Réalisation topologique d'un graphe}

\struck{$X(S, A, R) = (S \amalg (\vec{A} \times I))/$}

a) Segment associé à un ens.\ $a$ de card.\ $2$

[\,\struck{\ill{}} $\operatorname{seg}(a) \subset I^a$ ($I = [0,1]$)
\[
  \operatorname{seg}(a) = \Bigl\{ \{t_x\}_{x \in a} \Bigm| \textstyle\sum t_x = 1 \Bigr\} ;
  \quad \text{on a } a \hookrightarrow \operatorname{seg}(a)
\]

b)
\[
  \operatorname{réal}(S, A, R) = \Bigl( S \amalg \coprod_{a \in A}
  \operatorname{seg} R\{a\} \Bigr) \Big/ \text{relation d'équiv.}
\]

\struck{Réalisation} [ On a $X = \coprod_{a \in A} \operatorname{seg} R\{a\}$,
on a
\begin{tikzcd}[row sep=small]
  \partial X \arrow[r] \arrow[d, hook] & S \\
  X &
\end{tikzcd}
\ill{} « \uncertain{recollement} » $X : S \dashv \partial X \to S$ ]

« Réalisation \uncertain{géom.}\ rectiligne » (d'un graphe fini..) dans un
vectoriel (de dim finie…) sur $\mathbf{R}$ : homéom.\ dans
\note{« géom. » est ajouté au-dessus de la ligne}

\page{4}
\[
  \operatorname{réal}(S, A, R) \hookrightarrow E
\]
telle que l'image d'une arête soit un \emph{segment}.

\emph{Prop} Pour qu'un graphe fini admette une réalisation
\uncertain{géom.}\ rectiligne (dans un vectoriel convenable sur $\mathbf{R}$)
il f.\ et il s.\ qu'il soit strict — et alors \uncertain{il en admet} une
réalisation rectiligne dans $\mathbf{R}^3$. De façon précise, si
$S \hookrightarrow E$ ($\simeq \mathbf{R}^3$, v.a.)\ est une immersion telle
que
\[
  \forall x, y \in S \text{ distincts},\ x \neq y, \quad \text{on ait} \quad
  \underbrace{\overline{xy}}_{\substack{\text{segm.} \\ \text{joignant} \\ x \text{ à } y}}
  \cap\, S = \{x, y\},
\]
\marginal{i.e.\ trois points jamais alignés}
et $\forall x, y, z, t$ distincts, $(\overline{xy}) \cap (\overline{zt}) =
\emptyset$, $\exists !$ réalisation rectiligne
$\operatorname{réal}(S, A, R) \hookrightarrow E$ qui prolonge…
\note{« ($\simeq \mathbf{R}^3$, v.a.) », « distincts », « i.e. trois points
jamais alignés » et la condition sur $x, y, z, t$ sont au crayon, d'une main
plus rapide ; dans la marge gauche, au crayon, un petit dessin : un segment
dont les deux extrémités sont marquées}

Graphe connexe $=$ réal.\ géom.\ connexe $=$ quasi-\uncertain{non}
\note{la fin de la ligne n'est pas sûre}

\struck{\ill{}} « sommes » de deux autres graphes $\neq \emptyset$

\emph{NB} Tout graphe est de façon unique réunion de graphes connexes
\uncertain{disjoints}
\note{« disjoints » est ajouté sous la ligne, rattaché par un trait}

[graphe somme : $S \amalg S'$, $A \amalg A'$, $R \amalg R'$ … ]

$\to$ Graphe (\uncertain{multi}) polygonal (ou \emph{contour}
(combinatoire))
\[
  \bigl[\, \forall s \in S, \ \operatorname{card} R\{s\} = 2 \,\bigr]
\]
\note{« multi » est entouré, au-dessus de « polygonal » ; une flèche le relie
au « NB » qui précède}

\uncertain{Dénombrement} des \struck{graphes} \struck{polygonaux}
contours (combinatoires).
\struck{\ill{}} les \uncertain{composantes} d'un \struck{graphe polygonal}
contour \struck{polygonales} \uncertain{sont des} contours.

Classification des \struck{graphes polygonaux} contours conn.\ $\neq \emptyset$.

\emph{Prop} Un \struck{graphe polyg.}\ contour $\neq \emptyset$ connexe
\uncertain{est connu} à iso.\ près quand on connaît $\operatorname{card} S$
($= \operatorname{card} A$), qui est un cardinal \struck{fini ou infini}
\[
  2 \leq c \leq \aleph_0
\]
(strict ssi $c \neq 2$ i.e.\ $c \geq 3$)
\note{cette dernière ligne est au crayon}

\page{5}
Ils correspondent aux sous-groupes $H$ de $\mathbf{Z}$, distincts de
$\mathbf{Z}$ ($c = \operatorname{card} \mathbf{Z}/H$)…

\note{dans la marge gauche, deux dessins : au crayon, un contour à deux
sommets et deux arêtes ; à l'encre, un pentagone}

\emph{Orientation d'un contour} : déf.

\emph{Prop.} Soit $X = (S, A, R)$ un contour \struck{\ill{}} connexe
$\neq \emptyset$.

a) Soit \struck{Un} $g$ un autom.\ $g$ de $X$. \emph{Conditions équiv.}
\struck{\ill{}} :
\begin{itemize}
  \item[1°)] \struck{\ill{}} $\forall s \in S$, $s$ est lié à $gs$
  \item[2°)] $\forall a \in A$, $ga$ est lié à $a$
\end{itemize}
\note{sous « $g$ », dans la marge, « d'orientation » souligné deux fois :
l'énoncé se lit « un autom. $g$ de $X$ d'orientation »}

\struck{\ill{}} \uncertain{Déf.} On dit que $g$ est une \struck{\ill{}}
\emph{d'orientation} de $X$.

b) Soit $\omega(X) \subset \operatorname{Aut}(X)$ l'ens.\ des \struck{\ill{}}
d'orientation de $X$ (si $X$ est strict i.e.\ si $\operatorname{card} S \neq 2$).
Alors $\operatorname{card} \omega(X) = 2$, \struck{\ill{}} et si
$g \in \omega(X)$, on a $g^{-1} \in \omega(X)$, $\omega(X) = \{g, g^{-1}\}$.
\note{« $g^{-1} \in \omega(X)$ » est ajouté au-dessus de la ligne}

c) \struck{Soit $\operatorname{Aut}^{+}(X) \subset \operatorname{Aut}(X)$ le
s-groupe de $\operatorname{Aut} X$ \ill{} par les \ill{} $\operatorname{Aut} X$
tels que $g$ opère trivialement sur $\omega(X)$ [i.e.]}
Considérons
\[
  \operatorname{Aut}(X) \to \operatorname{Aut}(\omega(X)) \simeq \mathbf{Z}/2\mathbf{Z}.
\]
C'est un homom.\ surjectif, i.e.\ le noyau $\operatorname{Aut}^{+}(X)$ est un
s-groupe d'indice $2$.

De plus, $\operatorname{Aut}^{+}(X)$ est un groupe cyclique engendré par
\uncertain{tt} $g \in \omega(X)$.

\emph{Corollaire}
\[
  \operatorname{Aut}(X) \simeq \mathbf{Z}/2\mathbf{Z} \cdot \mathbf{Z}/H
  \quad \text{(produits ½ directs),}
\]
où $H$ $\to$ le ss-groupe de $\mathbf{Z}$ formé des $n \in \mathbf{Z}$ tels que
$g^n = 1$ ($g$ dans une orientation fixée)

\page{6}
\emph{Contour} \emph{orienté} « revient au $\equiv$ » que de se
donner un $n \in \mathbf{N}$, $n \geq 2$, et un \uncertain{tel} torseur sous
$\mathbf{Z}/n\mathbf{Z}$ (savoir $S$, --
\note{« tel » est ajouté au-dessus de la ligne}

$A$ -- \uncertain{choix} -- NB $\sigma : S \to A$ est compatible avec l'action
de $G^0 = \mathbf{Z}/n\mathbf{Z}$ i.e.\ avec l'action de $u_\omega$.

Se donner un contour \emph{connexe} équivaut :

la donnée de
\begin{itemize}
  \item[1°)] un groupe cyclique $G^0$ (d'ordre $n \geq 2$)
  \item[2°)] une \struck{\ill{}} partie de $G^0$ \struck{formée} de $2$ \ill{}
    de la forme $\{g, g^{-1}\}$, où $g$ est un \emph{générateur} (des
    parties de card $2$ ssi $n \neq 2$)
  \item[3°)] Un torseur $S$ sous $G^0$
\end{itemize}

\begin{itemize}
  \item[$\bullet$] On suppose si $n \geq 3$ \struck{\ill{}} ; pour une
    description
  \item[$\bullet$] générale, il faut se donner, au lieu de 2°), un
\end{itemize}
\begin{itemize}
  \item[2 bis)] un ens.\ $\Omega$ à deux éléments, et une application
    $\Omega \xrightarrow{i} G$ telle que, si $\Omega = \{a, b\}$, on ait
    $i(a)\, i(b) = 1$, $i(a)$ gén.\ de $G$
\end{itemize}

Contours orientés (\uncertain{non} \uncertain{connexes}) $=$ ensembles sur
lesquels $\mathbf{Z}$ opère (i.e.\ \uncertain{avec} \uncertain{autom.} \ill{})

Contours (non \uncertain{connexes}) $= \dots$

\page{7}
\emph{Orientation d'un contour.}

Choix pour toute arête $a \in A$ d'une « orientation de $a$ » i.e.\ d'une
origine $\sigma(a)$ de $a$, de telle façon que l'on ait

(O) \uncertain{Si} \struck{$R\{s\} = \{a, b\}$, alors $\exists$}
$s \in S$, $\exists !\, a \in R\{s\}$ tel que $s = \sigma(a)$
\marginal{$A \xrightarrow{\sigma} S$, $\forall a \in A$, $(\sigma(a), a) \in R$}\note{en
diagonale dans la marge gauche, avec quelques mots biffés illisibles ; un
trait sépare cette note du texte}

Réduction au cas connexe

\emph{Prop} Si $(S, A, R)$ est un contour connexe, il a exactement
\emph{deux} orientations ($\omega$ et $-\omega$).

\note{dans la marge gauche, deux petits dessins : deux arcs orientés mis bout
à bout, puis un arc de $s$ vers $u_\omega(s)$ marqué $\sigma^{-1}s$ ; plus
bas, l'arc $a$ d'origine $\sigma a$ suivi de l'arc $u_\omega(a)$, d'origine
$\sigma u_\omega(a)$}

\struck{Prop} \emph{Prop} Soit $\omega$ orientation des contours
$(S, A, R)$.

$\forall s \in S$, soit \struck{\ill{}} $u_\omega(s)$ l'extrémité de l'unique
arête \struck{\ill{}} $\sigma^{-1}(s)$ qui ait $s$ comme origine.
\[
  \bigl[ (u_\omega(s), \sigma^{-1}(s)) \in R,\ u_\omega(s) \neq s \bigr]
\]
$\forall a \in A$, soit $u_\omega(a)$ l'unique arc de $S$ dont l'origine soit
l'extrémité de $a$
\[
  \bigl[ (\sigma u_\omega(a), a) \in R,\ \sigma u_\omega(a) \neq \sigma a
  \ \text{ i.e. } u_\omega(a) \neq a \bigr]
\]

Alors $u_\omega : (S, A) \to (S, A)$ est un automorphisme du graphe
$(S, A, R)$. On a $u_{-\omega} = u_\omega^{-1}$. Les autom.\ $u_\omega$,
$u_\omega^{-1}$ sont caractérisés par les \struck{faits que} conditions
équivalentes
\begin{itemize}
  \item[1°)] $\forall s \in S$, $s$ lié à $gs$
  \item[2°)] $\forall a \in A$, $a$ lié à $ga$
\end{itemize}

Si $n = \operatorname{card} S = \operatorname{card} A$, l'ordre de $u_\omega$,
$u_\omega^{-1} = u_{-\omega}$ est $n$. Donc $u_\omega \neq u_{-\omega}$ ssi
$n \neq 2$ i.e.\ $n \geq 3$ ; donc $\omega \mapsto u_\omega$ est bijection de
$\operatorname{Or}(S, A, R)$ avec l'ens.\ des
$\{ g \in \operatorname{Aut}(S, A, R)$ satisfaisant les conditions équiv.\
1°, 2°$\}$, et le groupe cyclique $\simeq \mathbf{Z}/n\mathbf{Z}$ engendré par
$u_\omega$ est le \struck{\ill{}} groupe $G^0$ des automorphismes orientés
(i.e.\ de $(S, A, R; \omega)$) de $(S, A, R, \omega)$. $S$ et $A$ sont des
\emph{torseurs} sous $G^0$.
\note{« groupe $G^0$ » est ajouté dans la marge ; « orientés » et « (i.e. de
$(S, A, R; \omega)$) » au-dessus de la ligne ; « sous » est repassé}

\struck{La donnée d'un}

Torseur sous $\mathbf{Z}/n\mathbf{Z}$ ($n \geq 1$) $=$ ensemble à $n$
éléments muni d'un ordre « circulaire ». Se donner un contour
\note{la phrase s'arrête au bas du feuillet ; la page suivante reprend
ailleurs}

\section*{Ensembles à permutation}

\page{8}
\note{pages 8 et 9 à l'encre noire}

1) \emph{Ens.\ avec permutation} $(E, u)$ $\Leftrightarrow$
$\mathbf{Z}$-ens.

Classification : famille de cardinaux associés aux ss-groupes de $\mathbf{Z}$
$=$ famille de cardinaux $(a_i)_{i \in \mathbf{N}}$
\[
  E \text{ fini} \Longleftrightarrow
  \underbrace{n_0 = 0}_{\substack{\text{orbites} \\ \text{finies}}}
  \text{ et }
  \underbrace{\text{les } n_i \text{ finis}}_{\substack{\text{nb fini d'orbites de} \\ \text{cardinal } i \in \mathbf{N} \\ \text{donné}}},
  \underbrace{\text{nuls sauf un nb fini}}_{\text{nb fini d'orbites}}
\]
\note{la famille est notée $(a_i)$ à la ligne précédente, $(n_i)$ ici}

$u$ transitif ($\Leftrightarrow$ $u$ définit sur $E$ un ordre circulaire)
$\Leftrightarrow$ tous les $n_i$ sauf un \uncertain{sont} nuls
\note{« sauf un » est suivi, sous la ligne, de deux mots dont le premier est
lu « \uncertain{exemplaire} » : \ill{}}

\emph{Cor} Les classes d'éléments de $\mathfrak{S}_n$ mod.\ conjugaison
correspondent 1-1 aux « partitions » de $n$, i.e.\ aux familles d'entiers
$(c_i)_{1 \leq i \leq n}$ tels que $\sum i c_i = n$
\[
  \begin{array}{ll|ll}
    \mathfrak{S}_2 : 2 \text{ classes} & 1, \sigma_{12} &
      \mathfrak{A}_2 : 1 \text{ classe} & 1 \\
    \mathfrak{S}_3 : 3 \text{ classes} & 1, \sigma_{12}, \sigma_{123} &
      \mathfrak{A}_3 : 3\ \text{—} & 1, \sigma_{123}, \sigma_{321} \\
    \mathfrak{S}_4 : 5 \text{ classes} & 1, \sigma_{12}, \sigma_{12}\sigma_{34}, \sigma_{123}, \sigma_{1234} &
      \mathfrak{A}_4 : & \ldots \\
    \mathfrak{S}_5 : \ ? & & &
  \end{array}
\]
\struck{$\mathfrak{A}_4$ : $\sigma_{12}\sigma_{34}$, $\sigma_{1234}$}
\note{la ligne de $\mathfrak{A}_4$ est raturée et entourée ; on la donne
hors du tableau, et la lecture de ce qu'elle portait n'est pas sûre}

2) \emph{Rectification} Les contours orientés correspondent aux ens.\ à
permutation $(E, u)$ tels que $u$ n'ait pas de \emph{pt fixe} (car ce ne
sont pas des polygones : 1 côté, ds les graphes !)

3) Produit ½ direct de groupes

donnée
\[
  1 \to Z \xrightarrow{i} G \underset{q}{\overset{p}{\rightleftarrows}} \Gamma \to 1
  \qquad pq = \mathrm{id}_\Gamma
\]
(suite exacte, avec $q$ un homom.\ de groupes)

revient à donner
\[
  1 \to Z \to G \to \Gamma \to 1, \qquad \Gamma' \subset G, \quad
  p|\Gamma' \simeq \text{iso}
\]

ou à donner groupes $Z$, $\Gamma$ avec op.\ de $\Gamma$ sur $Z$
\[
  \Gamma \to \operatorname{Aut}_{\mathrm{gr}}(Z)
\]
loi de composition sur $G = Z \times \Gamma$

\struck{$(\gamma, z)(\gamma', z') = (\gamma\gamma', \operatorname{int}(\gamma'^{-1})(z).z')$}
\[
  (z, \gamma)(z', \gamma') = (z \operatorname{int}(\gamma)(z'), \gamma\gamma')
\]

\marginal{Ss-groupes $G'$ de $G$ avec $G' \supset Z$ $\Leftrightarrow$
ss-groupe $\Gamma'$ de $\Gamma$, \struck{\ill{}} et alors
$G' \simeq Z.\Gamma'$ (½ direct)}\note{en bas à droite, séparé du texte par un
trait oblique}

\marginal{Voir \uncertain{théorie} d'\ill{} \uncertain{extension} \ill{}
\ill{} $H^2(\Gamma, Z) \simeq$ \ill{}}\note{en bas à gauche, écrit en oblique
et en petits caractères, séparé du texte par un trait ; lecture très
incomplète}

\page{9}
\emph{Ex} Groupe \uncertain{affine} \struck{\ill{}} :
$\operatorname{Aff}(Z) \subset \mathfrak{S}_{Z}$ : un groupe $Z$
\note{« Ex » est relié par un trait à « Produit ½ direct » de la page
précédente}
\[
  1 \to Z \xrightarrow{i} \operatorname{Aut}_{\substack{\text{ens. à} \\ \text{gpe d'op.}}}(Z, Z_d)
  \underset{q}{\overset{p}{\rightleftarrows}} \operatorname{Aut}(Z) \to 1,
  \qquad \operatorname{Aut}(Z, Z_d) \hookrightarrow \operatorname{Aut}(Z) \times \mathfrak{S}_{Z_d}
\]
\note{au-dessus de $\operatorname{Aut}(Z, Z_d)$, relié par « $=$ », le nom
$\operatorname{Aff}(Z)$}
\[
  \gamma \in \operatorname{Aut}(Z), \qquad q(\gamma) = (\gamma, \gamma), \qquad
  i(g) = (\mathrm{id}_Z, \tau_g)
\]
\note{au-dessus de $\tau_g$, un « s » biffé}
\[
  \boxed{\operatorname{int}(q(\gamma)) | Z = \gamma}
  \qquad\qquad
  (\gamma, \rho) \in \operatorname{Aut}(Z, Z_d) \Longleftrightarrow
  \rho(gx) = \gamma(g)\rho(x) \quad \forall g, x \in Z
\]
\[
  \operatorname{Aff}(Z) \simeq Z.\operatorname{Aut}(Z) \quad
  \text{(produits ½ direct)}
\]
\note{après $\operatorname{Aff}(Z)$, un mot raturé ; sous lui, un autre,
également raturé}

Si $\Gamma \subset \operatorname{Aut}(Z)$ ($\Gamma$ ss-groupe) on a un
ss-groupe évident \struck{\ill{}} de $\operatorname{Aff}(Z)$, \uncertain{noté}
\[
  \operatorname{Aff}^{\Gamma}(Z) \simeq Z.\Gamma .
\]
\note{le mot raturé avant « de » porte en dessous « gr. d'op. »}

$\Bigl[$ Hom d'ens.\ à op.\ avec groupe variable
\[
  (E, G) \xrightarrow{(\rho, \gamma)} (E', G')
\]
\emph{NB} Si $E$ est un torseur sous $G$, et $E \neq \emptyset$, alors
$\rho$ détermine $\gamma$ car pour \struck{si} $x \in E$, $\forall g \in G$,
$\gamma(g)$ est déterminé par la condition
$\rho(g.x) = \gamma(g).\rho(x)$ $\Bigr]$
\note{« et $E \neq \emptyset$ » est ajouté au-dessus de la ligne ; « NB » est
précédé d'un premier « NB » biffé}

Cas $Z$ groupe sous-jacent à un vectoriel $V$, sur un corps $R$ (réels p.\ ex)
\note{« groupe sous-jacent à un » est ajouté au-dessus de la ligne}
\[
  \Gamma = \operatorname{Aut}_R(V) \subset \operatorname{Aut}_{\mathrm{gr}}(V)
\]
\[
  \operatorname{Aff}_R(V) \simeq V.\operatorname{Aut}_R(V)
\]
\[
  (\underset{V}{z}, \underset{\operatorname{Aut}_R(V)}{\gamma}) . (z', \gamma')
  = (z + \gamma(z'), \gamma\gamma')
\]

Si $V$ muni d'une forme quadratique \uncertain{non dégénérée} $Q$ $\to$ le
groupe des déplacements (transf.\ affines qui conservent la distance)
\[
  \operatorname{Dépl}(V, Q) \simeq V.\operatorname{Orth}(V, Q)
\]
\note{en indice de « Dépl », un « Aff » biffé : \struck{Aff}}
\note{au-dessus de « non dégénérée », un $Q$ ; la lecture de ces deux mots
n'est pas sûre}

\page{10}
$T$ torseur \uncertain{à droite} sous groupe $Z$

\struck{$\operatorname{Aut}(Z, T) \to \operatorname{Aut}(G)$}\note{premier
essai biffé, avec deux mots ajoutés dessous, également biffés}

\begin{tikzcd}
  1 \arrow[r] & Z \arrow[r, "i"] & \operatorname{Aut}(Z, T) \arrow[r, "p"] &
  \operatorname{Aut}(Z) \arrow[r] \arrow[l, dashed, bend left=30, "q_t"] & 1
\end{tikzcd}

$\forall t \in T$, on a une section $q_t$, définie par
\[
  \begin{cases}
    q_t(\gamma)(t) = t \\
    p(q_t(\gamma)) = \gamma
  \end{cases}
\]

Quelle est la relation entre $q_t$ et $q_{t'}$ ?

On aura $t' = t.g$, et on trouve
\[
  q_{t'} = \operatorname{int}(i(g)) \circ q_t : \gamma \mapsto
  i(g)\, q_t(\gamma)\, i(g)^{-1}
\]
$\Bigl[$ Pour le vérifier, il suffit de voir que
\[
  \underbrace{\operatorname{int}(i(g))(q_t(\gamma))}_{=\, i(g) q_t(\gamma) i(g)^{-1}}
  : t' \mapsto t'
\]
; or $t' = tg$

\struck{$q_t(\gamma)(t') = q_t(\gamma)(t)\,\gamma(g) = t.\gamma(g)$}

\struck{$\operatorname{int}(i(g))\, q_t(\gamma)(t') =$}
\begin{gather*}
  i(g)^{-1}(t') = t \qquad q_t(\gamma)\, i(g)^{-1}(t') = q_t(\gamma)\, t = t \\
  i(g)\, q_t(\gamma)\, i(g)^{-1}(t') = tg = t' \quad \text{ok} \ \Bigr]
\end{gather*}

Idem pour $\operatorname{Aut}^{\Gamma}(Z, T)$
\note{relié par un trait à la section $q_t$ du diagramme}

\emph{Ex} $Z = \mathbf{Z}_n$
\[
  \begin{cases}
    \operatorname{End} Z = \mathbf{Z}_n \\
    \operatorname{Aut} Z = \mathbf{Z}_n^{*} \supset \{\pm 1\} = \Gamma
  \end{cases}
  \quad (\text{\underline{NB} si } n = 2,\ +1 = -1 \text{ dans } \mathbf{Z}_n^{*} \ldots)
\]
\struck{\ill{}}
\[
  \begin{array}{l}
    \operatorname{Aff}(\mathbf{Z}_n) = \mathbf{Z}_n . \mathbf{Z}_n^{*}
      \quad \text{(produit ½ direct)} \\
    \quad \cup \\
    \operatorname{Aff}^{\{\pm 1\}}(\mathbf{Z}_n) = \mathbf{Z}_n . \{\pm 1\}
  \end{array}
\]
groupe diédral \uncertain{d'indice} $n$, $\mathbf{D}_n$ (on suppose
$n \neq 2$…)

Soit $T$ torseur sous $Z$ groupe cyclique d'ordre $n \geq 3$ muni d'une
famille $\{u, u^{-1}\}$ de $2$ gén.\ \struck{d'ordre $n$} inverses l'un de
l'autre.

\begin{tikzcd}
  1 \arrow[r] & Z \arrow[r, "i"] & \operatorname{Aut}^{\pm 1}(T, Z) \arrow[r, "p"] &
  \{\pm 1\} \arrow[r] \arrow[l, bend left=30, "q"] & 1
\end{tikzcd}

\page{11}
Données de $q$ équivalentes : la donnée de
$\sigma \in \operatorname{Aut}^{\pm 1}(Z, T)$
\[
  \begin{cases}
    p(\sigma) = -1 \\
    \sigma^2 = 1
  \end{cases}
\]
\[
  q_t(-1) = \sigma_t
\]
\[
  \sigma_{tg} = g \sigma_t g^{-1} = g \underbrace{\sigma_t g^{-1} \sigma_t^{-1}}_{g}
  \sigma_t = g^2 \sigma_t = \sigma_t g^{-2}
\]
\[
  \boxed{\sigma_{tg} = g^2 \sigma_t = \sigma_t g^{-2}}
\]

\emph{Corollaire} Si $n$ impair, $t \mapsto \sigma_t$ est une bijection
de $T$ avec l'un \struck{des} $p^{-1}\{-1\}$ ; donc si on a un nb impair de
côtés, l'application $t \mapsto \sigma_t$ qui, à \uncertain{un} sommet du
polygone, associe l'unique autom.\ non orienté qui fixe ce sommet, est
bijection.

\emph{NB} Si $n$ pair, elle n'est pas bijective, \struck{mais} ses
fibres sont toutes de cardinal $2$…

\struck{\ill{}} Sommet \emph{opposé}, symétrie canonique

\marginal{$\{u, \bar{u}\} \subset Z$ gén.\ mais d'ordre \ill{}, \ill{}
$u \bar{u} = s$, \uncertain{cela a lieu} \ill{} \ill{} $\sigma_t = \sigma_{t'}$
ssi $t, t'$ \ill{} $\{1, s\}$}\note{en diagonale dans la marge gauche, dans un
cadre relié à « Parenthèse » ; lecture très incomplète}

\emph{Parenthèse} Produits fibrés d'un

\begin{tikzcd}
  P \arrow[r, dashed, "p'"] \arrow[d, dashed, "q'"'] & F \arrow[d, "q"] \\
  E \arrow[r, "p"] & S
\end{tikzcd}
\note{dessiné avec $E \to S$ en haut, $F$ en bas, et les flèches $q$, $q'$
montantes ; redessiné ici en carré}
\[
  P \subset E \times F, \qquad P = \{ (x, y) \mid px = qy \}
\]

Solution d'un pb universel dans la cat.\ des ens :

Si $E$, $F$, $S$ des groupes, $p$, $q$ hom.\ de groupes, alors $P$ un groupe
(et solution du pb univ.\ analogue dans cat.\ des groupes)
\[
  \operatorname{Ker} p \simeq \operatorname{Ker} p', \qquad
  \operatorname{Ker} q \simeq \operatorname{Ker} q'
\]
Si $p : E \to S$ surjectif (bijectif) alors $P \to F$ surjectif (bijectif).

[Langage des \struck{diagrammes} carrés cartésiens]

Image inverse d'ex.\ de groupes

\page{12}
\begin{tikzcd}[row sep=small]
  1 \arrow[r] & Z \arrow[r] & G \arrow[r] & \Gamma \arrow[r] & 1 & (E) \\
  1 \arrow[r] & Z \arrow[r] \arrow[u, "\mathrm{id}"'] & G' \arrow[r] \arrow[u] &
  \Gamma' \arrow[r] \arrow[u] & 1 & (E')
\end{tikzcd}

Si l'extension $(E)$ splittée, l'ext.\ $E'$ aussi.

\emph{Remarque} Toute extension splittée $G$ de $\Gamma$ par $Z$ provient
par image inverse d'une ext.\ splittée $Z.\operatorname{Aut}(Z)$.

4) \uncertain{Ens.}\ de petit cardinal

a) \emph{Cardinal 2} \quad Ens de cardinal $2 \longleftrightarrow$ torseurs
sous $\mathbf{Z}_2$ (ou encore, \emph{espace affine sur $\mathbf{F}_2$}
\struck{\ill{}} ou encore, espace affine de dim $1$ sur corps $\mathbf{F}_2$ ---
car les espaces vect.\ de dim $1$ sur $\mathbf{F}_2$ sont can.\ isomorphes
…)
\[
  \boxed{\operatorname{Aff}(1, \mathbf{F}_2) \simeq \operatorname{Aff}(\mathbf{Z}_2)
  \simeq \mathfrak{S}_2}
  \qquad (\mathfrak{A}_2 = \{1\})
\]

b) \emph{Cardinal 3}
\[
  \mathfrak{S}_3 \simeq
  \underset{\substack{\wr \\ \mathbf{F}_3^{+}.\mathbf{F}_3^{*} \\
  \mathbf{Z}_3 \quad \mathbf{Z}_2}}{\operatorname{Aff}(1, \mathbf{F}_3)}
  \simeq \mathrm{Gl}(2, \mathbf{F}_2) \simeq
  \underset{\substack{\| \\ \mathrm{Gl}(2, \mathbf{F}_2)/\mathbf{F}_2^{*}}}{\mathrm{GP}(1, \mathbf{F}_2)}
\]
\note{sous $\mathbf{F}_3^{+}$ et $\mathbf{F}_3^{*}$, des signes $\simeq$ et $=$
verticaux les relient à $\mathbf{Z}_3$ et $\mathbf{Z}_2$}
\[
  \mathfrak{A}_3 \simeq \mathbf{F}_3^{+} \ (\simeq \mathbf{Z}_3)
\]

\page{13}
c)
\begin{align*}
  \mathfrak{S}_4 &\simeq \operatorname{Aff}(2, \mathbf{F}_2) \simeq \mathrm{GP}(1, \mathbf{F}_3) \\
  \mathfrak{A}_4 &\simeq \operatorname{Aff}(1, \mathbf{F}_4) \quad ?
\end{align*}

d)
\begin{align*}
  \mathfrak{S}_5 &\simeq \mathrm{GP}(1, \mathbf{F}_5) \\
  \mathfrak{A}_5 &\simeq \mathrm{GP}(1, \mathbf{F}_4)
\end{align*}

Quelques groupes classiques sur un corps $k$
\[
  \left.
  \begin{array}{l}
    1 \to k^{*} \to \mathrm{Gl}(V) \to \mathrm{GP}(V) \to 1 \\
    1 \to k^{*} \to \mathrm{Gl}(n, k) \to \mathrm{GP}(n-1, k) \to 1
  \end{array}
  \right\} \text{non splittées en général}
\]

\marginal{$\mathbf{P}(V) \simeq V/k^{*}$ : $\mathrm{GP}(V)$ \uncertain{y
opère}}\note{dans la marge gauche, entouré}

Structure de $k$-vectoriel (\uncertain{sur un} $V$ de dim $n$)
$\longleftrightarrow$ donnée d'un élément de
\[
  \operatorname{Isom}_{\text{ens}}(k^n, V)/\mathrm{Gl}(n, k)
\]
\note{« de dim $n$ » est ajouté au-dessus de la ligne}

Structure de $k$-espace affine de dim.\ $n$ sur \struck{$E$} un $E$
$\longleftrightarrow$ donnée d'un élément de
$\operatorname{Isom}_{\text{ens}}(k^n, V)/\operatorname{Aff}(n, k)$

$\Bigl[$ De façon générale, \struck{si $Z$ est un groupe}, la notion de
(\uncertain{torseur} \uncertain{Ed} \uncertain{abstrait}, disons) $E$
\struck{espace aff.}\ sous un groupe non précisé que l'on considère comme
structure sur le seul ens.\ $E$, \struck{\ill{}} \uncertain{précisément} donnée
d'un \uncertain{\ill{}} $E$ \struck{\ill{}} \struck{$(x, y, z) \mapsto$} \ill{}
couple $\Gamma \subset \mathfrak{S}_E$ ss-groupe, tel que …\ ;
\uncertain{muni} par ens.\ à groupe d'op.\ \emph{fidèle}…

Une structure d'espace \emph{affine sur $k$} est alors la donnée d'un tel
$\Gamma$, avec \struck{\ill{}} commutatif, et un hom.\ d'anneaux unit.\
\[
  k \to \operatorname{End}_{\mathbf{Z}}(\Gamma) \quad \ldots \ \Bigr]
\]
\note{passage rapide et en partie biffé ; la lecture des mots entre « torseur »
et « couple » n'est pas sûre}

Structure de $k$-espace projectif de dim.\ $n$ sur $P$
$\overset{\text{(déf)}}{\longleftrightarrow}$ donnée d'un élément de
\[
  \operatorname{Isom}_{\text{ens}}(\mathbf{P}^{n-1}(k), P)/\mathrm{GP}(n-1, k)
  \qquad \mathbf{P}^{n-1}(k) = (k^n - 0)/k^{*}
\]
\note{la dimension notée $n$ dans la phrase et $n-1$ dans la formule, comme
sur la page}
\[
  \mathrm{Gl}(n, k) \longrightarrow \operatorname{Aff}(n, k) \longrightarrow \mathrm{GP}(n, k)
\]
vect.\ de dim $n$ sur $k$ $\longmapsto$ (oubli origine) espace affine de dim $n$
sur $k$ $\longmapsto$ espace projectif associé à l'esp.\ vect.\ enveloppant
(esp.\ proj.\ de dim $n$ sur $k$)
\note{sur la page, les trois correspondances sont écrites en colonnes sous
les trois groupes, les mots « oubli origine » et « associé à l'esp. vect.
enveloppant » en petits caractères sous les flèches}

\page{14}
\emph{Exercice} Réalisation régulière d'un polygone.

Soit $E$ plan affine euclidien sur $\mathbf{R}$ (i.e.\ muni d'une dist…\
sur le plan $V$ des translations). Une réalisation rectiligne d'un polygone
$(S, A, R)$ est dite \emph{régulière} si tous les côtés du polygone ont
même longueur, \uncertain{et tous les angles itou},
\struck{\ill{}} \uncertain{et} unimodulaire si la dite longueur est $1$.
\note{« et tous les angles itou » est ajouté dans une bulle au-dessus de la
ligne}

a) Existence, grâce à équation cyclotomique $z^n = 1$ dans $\mathbf{C}$.

b) Prouver que si $(V \supset P \supset S)$ et $(V' \supset P' \supset S')$
sont deux $n$-polygones ($n \geq 3$) munis d'une réal.\ rectiligne rég.\
(\uncertain{supposés} unimodulaires) il existe une \struck{\ill{}} similitude
(\uncertain{et} \uncertain{une} \uncertain{seule}) \struck{unique} de $V$ sur
$V'$, qui transforme $P$ en $P'$, $S$ en $S'$ ; \struck{\ill{}} plus précisément,
pour \uncertain{la} \uncertain{donnée} (\uncertain{combinatoire}) de
$(S, A, R)$ sur $(S', A', R')$, $\exists !$ similitude \uncertain{(isométrie)}
de $V$ sur $V'$ qui l'induise. \struck{\ill{}}
\note{« (isométrie) » est ajouté au-dessus de « de $V$ »}

\uncertain{c)} \emph{Cas} Notion d'enveloppe affine (ou convexe, de l'ens.\
des sommets) \uncertain{et} \uncertain{centre} de gravité d'un polygone
combinatoire.
\note{« de l'ens. des sommets » est ajouté sous la ligne ; la lettre de
l'item est surchargée}

Si Orientations des polygones $\longleftrightarrow$ orientation de son
enveloppe

\page{15}
Polygone régulier (ou réalisation régulière d'un polygone combinatoire) :
définitions possibles (côtés égaux, angles égaux, \struck{\ill{}} les deux,
transitivité groupe des similitudes…) Existence, unicité modulo
homothétie : tout isomorphisme combinatoire est induit par une unique
similitude.
\note{« les deux, » est ajouté au-dessus de la ligne, un mot biffé devant ;
dans la marge gauche, un petit dessin : deux segments formant un angle}

Notion d'enveloppe affine (ou vectorielle) d'un polygone combinatoire…\
\uncertain{sans} précision

\emph{Fonctorialité de $X \times_G Y$ par : $(G, X, Y)$…}
\[
  P \times_G E = {}^{P}E = F \qquad
  \begin{array}{l}
    P \ G\text{-torseur} \\
    E \ G\text{-ensemble}
  \end{array}
\]
\[
  P \to \operatorname{Bij}(E, F) \quad \text{compatible aux opérations de } G
  \text{ : droite}
\]
\emph{NB} Si $G$ opère fidèlement sur $E$, $G$ opère « sans pts fixes »
sur $\operatorname{Bij}(E, F)$, i.e.\ $u.g = u$ ($g \in G$,
$u \in \operatorname{Bij}(E, F)$) implique $g = 1$, on a donc les
$G_u$ ($u \in B$), \ill{} : $\{1\}$, \ill{} $B$ est une \ill{} de
torseurs…\ Dans ce cas (et dans ce cas seulement)
$P \to \operatorname{Bij}(E, F)$ est injectif
\note{« et dans ce cas seulement » est ajouté sous la ligne}
\[
  P \hookrightarrow \operatorname{Bij}(E, F),
\]
$F = {}^{P}E$ \uncertain{est} muni de 2 structures
\note{au-dessus de « $F =$ », un « $P_E$ » relié par un trait}
$\Bigl[$ $E$ -- dans les cas, $F$ \ill{} : donnée d'une orbite de $G$ dans
$\operatorname{Bij}(E, F)$. [On \uncertain{verra} que (si $G$ opère fid.\ sur
$E$) la donnée \struck{d'une $(G, E)$ \uncertain{structure}} \ill{} d'une
$(G, E)$-structure \ill{} équivaut à celle d'un \emph{torseur} sous
\uncertain{$G$} \ill{} \ill{}.] $\Bigr]$
\marginal{$(G, E)$-structure sur un $F$}\note{dans la marge gauche, en
oblique, relié par une accolade au crochet ; la fin du paragraphe, serrée en
petits caractères au bas du feuillet, n'est lue qu'en partie}

\page{16}
sur $E$ invariante par $G$ se transporte sur $F$.
(Ex : structures de groupe, d'anneau, de module sur un anneau fixé, d'espace
top.\ etc…)

Cas de plusieurs ens.\ de base $(E_i)_{i \in I}$, \struck{\ill{}} munis d'une
structure mixte (p.\ ex.\ module $E$ sur anneau $A$, le couple $(A, E)$)
alors $({}^{P}E_i)_{i \in I}$ sont munis de la structure correspondante.

\struck{\ill{}}

\emph{Applications} \struck{\ill{}} (1) Yoga : structures isomorphes : \uncertain{les}
(\uncertain{str}) données $E \longleftrightarrow$ torseurs sous
$G = \operatorname{Aut}(E)$
\note{« str » est ajouté au-dessus de la ligne}
\[
  \begin{array}{c}
    F \longmapsto \operatorname{Isom}_{\mathcal{C}}(E, F) = P \\
    P_E = P \times_G E \longleftarrow P
  \end{array}
\]
\[
  F \longmapsto \underset{\operatorname{Isom}_{\mathcal{C}}(E, F)}{P}
  \longmapsto P \times_G E \overset{?}{\xrightarrow{\ \sim\ }} F
\]
\[
  \begin{array}{ll}
    P \times E \longrightarrow F & \\
    \operatorname{Isom}(E, F) \times E \to F & (u, x) \mapsto u(x) \\
    & (ug)(x) = u(gx) \quad \text{ok}
  \end{array}
\]
\[
  P \to F = P \times_G E \longmapsto \operatorname{Isom}_{\mathcal{C}}(E, F)
  \overset{\sim}{\longleftarrow} P
\]
\[
  i_u : x \mapsto u * x \longleftrightarrow u \qquad
  i_{ug} = i_u . g \quad \text{i.e.}\quad
  ug * x = u * gx \quad \text{ok}
\]
\note{devant $i_{ug}$, un premier essai biffé : \struck{$i_{ug}$}}

\page{17}
(2) Donnée d'un $G$-torseur $T$ $\longleftrightarrow$ donnée d'un
ensemble $F$ muni d'une orbite de $G$ opérant sur $\operatorname{Bij}(E, F)$
[quand $E$, \struck{\ill{} fid.\ \ill{} $\operatorname{Bij}(E, F) \to$}
$G$-ens.\ fidèle, donné une fois pour toutes] C'est un cas particulier de
(1) car $G = \operatorname{Aut}_{\mathcal{C}}(E)$, $\mathcal{C}$
étant l'espèce de $\mathcal{C}$-structures du type $(G, E)$.
\marginal{i.e.\ $F$ muni d'une $(G, E)$ structure}\note{dans la marge gauche,
en oblique, relié par un trait à l'accolade}

(3) Soit $\mathcal{C}$ une espèce de structures (sur un ens.\ de
base, disons), $E$ une structure de type $\mathcal{C}$,
$G = \operatorname{Aut}_{\mathcal{C}}(E)$. Alors, pour un ens.\ $F$, la donnée
d'une structure de type $\mathcal{C}$ \emph{isom}\add{orphe} à $E$
\emph{équivaut} à celle d'une $(G, E)$-structure. De cette façon,
(1) devient un particulier de (2).

(4) Enveloppe vectorielle d'un polygone combinatoire.

On considère le polygone \struck{\ill{}} combinatoire standard
\struck{$[1, n]$} $P_n$, sa réalisation (rectiligne) régulière
\uncertain{Arcs}
\[
  i \mapsto \exp \frac{2 i \pi}{n} = \zeta^i
\]
(il faut prendre $n \geq 3$ pour que ceci soit une réalisation rectiligne…\
!)
\note{« (rectiligne) » est ajouté au-dessus de « régulière » ; « $n \geq 3$ »
sous la ligne}

On note que $D_n = \operatorname{Aut}(P_n)$ opère sur $\mathbf{C}$ ---
$E =$ plan euclidien $\mathbf{C}$ (muni de $|z|$) ---
par
\[
  \begin{cases}
    \tau_i \longmapsto \text{homothétie par } \zeta^i \\
    \sigma \longmapsto z \mapsto \bar{z}
  \end{cases}
\]

\page{18}
La donnée d'un polygone combinatoire $P$ à isomorphisme près équivaut à la
donnée d'un $D_n$-torseur $T$ ($P = {}^{T}P_n$), et : $T$ \struck{\ill{}}
associé ${}^{T}E$.

L'immersion $S(P_n) \hookrightarrow E$ (\uncertain{liée} aux $D_n$-opérations)
définit une immersion
\[
  S(P) = {}^{T}S(P_n) \hookrightarrow {}^{T}E = \operatorname{Env}(P)
\]
\note{« liée » n'est pas sûr ; la marge de cette ligne est serrée}

\[
  \underline{\mathrm{Hom}}(\mathbf{e}, G) \overset{\approx}{\longrightarrow}
  \underline{\mathrm{Tors}}(G)
\]
\emph{NB}
\[
  \underline{\mathrm{Hom}}(\underline{\mathrm{Tors}}(G), \mathcal{C})
  \overset{\approx}{\longrightarrow} \text{catégorie}\ G\text{-}(\mathcal{C})
\]
des objets $X$ de $\mathcal{C}$ munis d'une opération à gauche de $G$,
$G \to \operatorname{Aut}(X)^{\circ}$
\note{« gauche » remplace « droite », biffé. « $G$-($\mathcal{C}$) » est
écrit au-dessus d'un premier nom biffé,
\struck{$\underline{\mathrm{Hom}}(G, \mathcal{C})$}}
\[
  \begin{array}{rcl}
    F & \longmapsto & F(G_d) \\
    (P \mapsto {}^{P}X) & \longleftarrow & X
  \end{array}
\]
\note{devant le premier $F$, un mot biffé : \struck{\ill{}}}

Soit $\mathcal{C}$ catégorie, $E \in \operatorname{Ob} \mathcal{C}$,
$\mathcal{C}_E$ = cat.\ des objets de $\mathcal{C}$ isom.\ à $E$, avec iso,
$G = \operatorname{Aut}_{\mathcal{C}}(E)$
\[
  \mathcal{C}_E \overset{\approx}{\longrightarrow} \underline{\mathrm{Tors}}(G)
\]

\begin{tikzcd}[column sep=large]
  \underline{\mathrm{Hom}}(\underline{\mathrm{Tors}}(G), \mathcal{E})
  \arrow[r, "\approx"] \arrow[d, "\approx"'] &
  G\text{-}(\mathcal{E}) \arrow[dl, "{\Phi,\ \approx}" description] \\
  \underline{\mathrm{Hom}}(\mathcal{C}_E, \mathcal{E}) &
\end{tikzcd}
\note{la flèche du haut porte au-dessous un nom biffé,
$\underline{\mathrm{Hom}}(G, \mathcal{E})$, remplacé par
« $G$-($\mathcal{E}$) » ; la flèche oblique, dessinée à double pointe, est
surchargée ; à côté, $\Phi \leftrightarrow \Phi(E)$, puis en pointillés
$X \mapsto (F \mapsto {}^{\operatorname{Isom}(E, F)}X)$}

Donc : la catégorie des foncteurs de $\mathcal{C}_E$ dans $\mathcal{E}$ est
équivalente (de façon can.\ : \uncertain{isom.}\ unique près) : celle des
$G$-objets de $\mathcal{E}$
\marginal{\emph{cf}\ le groupoïde connexe, $E \in \operatorname{Ob}
\mathcal{C}_0$, $G = \operatorname{Aut}(E_0)$ :
$\underline{\mathrm{Hom}}(\mathcal{C}_0, \mathcal{E}) \to G\text{-}(\mathcal{E})$,
$\Phi \mapsto \Phi(E)$ équivalence de catégories.}\note{dans la marge gauche,
en oblique, à gauche d'un long trait vertical qui porte « DEA » et embrasse
tout le bas de la page}

\page{19}
Supposons que : $\forall$ polygone comb.\ de type $P$ (disons) \uncertain{on}
veuille faire correspondre un $\mathbf{R}$-vectoriel \struck{\ill{}} (disons),
« de façon compatible avec le transport de structure », \struck{Alors} $E(P)$.
Alors $E(P_n) = E$ est un vectoriel, avec $D_n = \operatorname{Aut}(P_n)$
opérant (\uncertain{à} \uncertain{g}.). Mais la donnée de $E$ avec cette
opération permet de reconstituer $E(P)$ pour tout $P$, comme étant
\struck{$E_T$} ${}^{T}E$, où $T = \operatorname{Isom}(P_n, P)$.
\note{« $= E$ » est ajouté au-dessus de $E(P_n)$ ; « les » au-dessus de
« Alors » biffé}

(5) \struck{Changement de} Extension du groupe d'opérateurs

$G \to G'$, $E$ $G$-ens.\ (à gauche), on définit
\[
  {}^{G'}E = E \times_G G'
  \quad \text{comme } G'\text{-ens.\ à g.}
\]
\note{devant ${}^{G'}E$ et devant $E \times_G G'$, une lettre biffée
chaque fois : \struck{\ill{}}}

\emph{NB} Si $F$ un $G'$-ens.\ à droite, dans
\[
  \left\{
  \begin{array}{l}
    G\text{ ens.\ à droite}, \ldots \\
    (E \times_G G') \times_{G'} F \simeq E \times_G (G' \times_{G'} F) \simeq E \times_G F
  \end{array}
  \right.
\]
\struck{$F \times_{G'} (G' \times_G E) \simeq (F \times_{G'} G') \times_G E \simeq F \times_G E$ i.e.}

\struck{$F \times_{G'} E' \simeq F \times_G E$}
\[
  E' \wedge_{G'} F \simeq E \wedge_G F
\]
Si $E = T$ torseur, $E' = T'$ \uncertain{aussi}, donc \ill{}.
\[
  {}^{T}F \simeq {}^{T'}F
\]
\note{au-dessus de « torseur », un signe biffé ; après ${}^{T}F$ et après
${}^{T'}F$, un exposant biffé chaque fois : \struck{\ill{}}}
i.e.\ \uncertain{Pour} tordre $F$ par \struck{\ill{}} $G$-torseur $T$, quand
$G$ opère

\page{20}
\emph{via} réduction d'un groupe $G$ ; il suffit de tordre par le
torseur ${}^{G'}T = T'$. (Ex : si $F$ muni de quelques structures,
\struck{diverses} \ill{} on prend $G' = \operatorname{Aut}(F)^{\circ}$…\
et $F^{T}$ \ill{} \ill{} que la structure \uncertain{déduite} sur
$G'$-torseur $T'$…)
\note{« quelques » est ajouté au-dessus de la ligne ; « déduite » n'est pas
sûr}

(6) Opérations \struck{sur les} du produit sur les torseurs sous
groupe \emph{commutatif}. Opération à droite \uncertain{produit}
\[
  X \times_G Y \times_G Z \ \ldots \ X \wedge Y \wedge Z
\]
produit associatif et commutatif, à isom.\ can.\ près ;
$\mathbb{1} = G_s$ pour la \uncertain{dite} \uncertain{mult.}
\note{devant $X \wedge Y \wedge Z$, un $X$ biffé ; « à isom. can. près » et
« pour la dite mult. » sont en petits caractères, en partie sous la ligne}
\[
  X \times_G Y \simeq (X \times Y) \times_{G \times G} G
  \qquad
  (X \times Y \times \ldots) \times_{G^n} G
\]
\note{sous $G^n$, un indice biffé}

\emph{Prop} Un produit de torseurs est un \emph{torseur}.

\struck{\ill{}}

\emph{Prop} $\forall$ torseur $X$, $\exists$ torseur \struck{$Y$} $Y$ et
iso.\ $X \wedge Y \simeq \underline{G}$. \uncertain{Pour} \uncertain{un} \uncertain{tel}
$Y'$, $\exists$ iso.\ unique $Y \simeq Y'$ tel que…
\[
  Y \simeq X \times_G (G \xrightarrow{g \mapsto -g} G)
\]

(7) \emph{Ex} \struck{Si Orientations, \uncertain{d'un} $X \wedge Y$
\ill{} \ill{} : \ill{}} \struck{\ill{} \ill{}}
\note{deux lignes biffées : au-dessus de la première, « $X \wedge Y$ » ;
de la seconde ne se lisent que quelques lettres}

\struck{\emph{Ex} Soit $B$ ens.\ fini avec \uncertain{rel.}\ \ill{}. $R$}
\struck{à relation des conditions 2, $B \to I = B/R$}
\struck{Considérons $\bigwedge_{i \in I} B_i$, un élément}
\note{les trois dernières lignes, encadrées, sont barrées de traits obliques ;
au-dessus de $\bigwedge$, un astérisque. La suite, section « 8 », ouvre le
lot 2 (page 21)}
\end{document}
